\subsection{A uniform genus mean for primitive sublattices}
\label{sec:primitive-sublattice-genus-mean}

We now supply the arithmetic mean that counts primitive sublattices
in the original fixed-label return argument.  The proof uses
the exact mass formulas of Section~\ref{sec:maximal-hermitian-mass},
finite discriminant modules, and the classical local-to-global
principle for Hermitian spaces.

Retain a CM extension \(E/F\), unramified at all finite places,
and write
\[
 d=[E:\mathbf Q]=2f,\qquad \ell=\log_2d.
\]
Let \(V\) be a positive definite Hermitian \(E\)-space of rank
\(j+s\), where \(j,s\geq1\).  For an integral lattice \(A\) in
any such space, write \(\Delta(A)=|A^\#/A|\).
If \(L\subset V\) is an integral lattice, define
\[
 \mathcal U_{j,n}(L)=
 \{U\subset V:\dim_EU=j,\ \Delta(L\cap U)=n\}.
 \tag{PG1}
\]
Here \(U\) is an \(E\)-linear subspace.  Its intersection
with \(L\) is saturated in \(L\), so these are precisely the
primitive sublattices of rank \(j\) and discriminant \(n\),
specified by their rational span.

\begin{theorem}[The primitive-sublattice mean]
\label{thm:primitive-sublattice-genus-mean}
Let \(n,N\geq1\) be integers and let \(\mathcal G\) be any
isometry-stable collection of integral lattices in \(V\) with
\(\Delta(L)=N\).  Then
\[
 \sum_{[L]\in\mathcal G}
        \frac{|\mathcal U_{j,n}(L)|}{|\operatorname{Aut}(L)|}
 \leq C(E,j)C(E,s)\,
      n^{j/2}(nN)^{(9s+8+4\ell)/2}.
 \tag{PG2}
\]
The set \(\mathcal U_{j,n}(L)\) is finite.  The sum and all its
automorphism weights are the original genus-counting ones.
The bound does not require the discriminant module \(L^\#/L\)
to be cyclic.
\end{theorem}

\begin{proof}
\smallskip\noindent\emph{The number of rational subspace types.}
Consider the positive definite Hermitian spaces of rank \(j\)
that contain an integral lattice of discriminant \(n\).
For any such space, choose a maximal integral lattice \(M\).
Its finite nonunimodular places form a set \(B\) of inert
primes of \(F\), and
\[
 \prod_{v\in B}q_v\leq\sqrt n.
 \tag{PG3}
\]
Indeed \(\Delta(M)=(\prod_{v\in B}q_v)^2\), and the
discriminant-index identity (MM15) applies to the lattice
of discriminant \(n\) contained in \(M\).

The set \(B\) determines the rational Hermitian space up to
isometry.  At an inert finite place the maximal local lattice
has determinant valuation zero or one according to membership
in \(B\), and its unit diagonal coefficients can be removed
by local norms.  At a split place the rational form is unique.
At infinity all forms are positive definite.  The Hermitian
local-to-global principle therefore applies
\cite[Theorem 2.4.1]{kirschmer-herm}.

Map \(B\) to the squarefree ideal \(\prod_{v\in B}v\) of
\(\mathcal O_F\).  By (PG3), its norm is at most \(n\).
The ideal count (MM17), used for \(F\) and bounded by
the larger degree \(d\), shows that there are at most
\[
 \sum_{a=1}^{n}a^\ell\leq n^{1+\ell}
 \tag{PG4}
\]
possibilities.  This counts all rational spaces under
consideration, whether or not they occur in \(V\).

Isometric subspaces of \(V\) of dimension \(j\) lie in
the same orbit of its unitary group.  One may use Witt's
extension theorem here.  Alternatively, their orthogonal
complements have the same local determinant parities and
the same signatures, hence are isometric by the same
local-to-global principle.  Combining the two isometries
in the orthogonal direct sums gives an isometry of \(V\).
Thus at most \(n^{1+\ell}\) subspace orbits occur.

\smallskip\noindent\emph{The discriminant module of an extension.}
Fix one representative \(U\), put \(W=U^\perp\), and let
\[
 A=L\cap U,\qquad B=L\cap W,\qquad b=\Delta(B).
\]
Both \(A\) and \(B\) are primitive in \(L\), and
\[
 A\perp B\subset L\subset A^\#\perp B^\#.
 \tag{PG5}
\]
Write \(G_A=A^\#/A\) and \(G_B=B^\#/B\), of cardinalities
\(n,b\).  The image of \(L/(A\perp B)\) in \(G_A\oplus G_B\)
projects injectively to either factor.  For example, if its
first projection is zero, subtracting an element of \(A\)
leaves an element of \(L\cap W=B\).  Hence that image is
the graph of an injective homomorphism between submodules
of \(G_A\) and \(G_B\).  If its order is \(t\), then
\[
 N=\frac{nb}{t^2},\qquad t\leq n,\qquad
                         1\leq b\leq nN.
 \tag{PG6}
\]

Every submodule of \(G_B\) can be generated by at most \(s\)
elements over \(\mathcal O_E\).  To verify this even when
the lattices are not free, localize at each prime.  The
module \(B^\#\) is locally free of rank \(s\), so every
submodule of its finite quotient by \(B\) has at most
\(s\) cyclic factors over the local discrete valuation
ring.  Choose local generators, pad the lists with zeros,
and combine the finitely many lists by the Chinese remainder
theorem.  This gives \(s\) global generators.

Consequently there are at most \(b^s\) possible submodules
of \(G_B\): count their ordered generating lists of length
\(s\).  For each one, there are at most \(n^s\) homomorphisms
to \(G_A\).  Therefore at most
\[
                            (nb)^s
 \tag{PG7}
\]
lattices \(L\) can occur for fixed \(A,B\).  Each graph
determines its inverse image in (PG5).  Requiring that the
graph be isotropic for the discriminant pairing, that both
projections be injective, and that its order give the
specified \(N\) only decreases this count.  No cyclicity
of \(L^\#/L\) has been used.

\smallskip\noindent\emph{The automorphism weights.}
The stabilizer of \(U\) in the unitary group of \(V\)
is the product of the unitary groups of \(U\) and \(W\).
For fixed isometry classes of \(A,B\), its remaining
action on the extension graphs is the finite action of
\(\operatorname{Aut}(A)\times\operatorname{Aut}(B)\).
The sum of \(1/|\operatorname{Aut}(L,U)|\) over its
graph orbits is exactly the number of graphs divided by
\(|\operatorname{Aut}(A)||\operatorname{Aut}(B)|\).
Thus (PG6)--(PG7) give, with \(X=nN\),
\[
 \begin{split}
 &\sum_{\substack{[(L,U)]\\
                  \Delta(L)=N,\ \Delta(L\cap U)=n}}
               \frac1{|\operatorname{Aut}(L,U)|}\\
 &\quad\leq
 \left(\sum_{\substack{[A]\text{ in }U\\\Delta(A)=n}}
                       \frac1{|\operatorname{Aut}(A)|}\right)
 \left(\sum_{b=1}^{X}
        \sum_{\substack{[B]\text{ in }W\\\Delta(B)=b}}
                       \frac1{|\operatorname{Aut}(B)|}\right)
                           (nX)^s .
 \end{split}
 \tag{PG8}
\]
The subspace representative \(U\) is fixed on this line.
Theorem~\ref{thm:integral-lattice-fixed-discriminant-mass}
bounds the two factors by
\[
 C(E,j)n^{(j+\ell)/2},\qquad
 C(E,s)X^{\,1+(s+\ell)/2},
 \tag{PG9}
\]
respectively.  These mass bounds hold for the actual two
rational spaces, independently of their local determinant
parities.

Finally sum (PG8) over at most \(n^{1+\ell}\) rational
subspace orbits.  Since \(n\leq X\), the result is at most
\[
 \begin{split}
 &C(E,j)C(E,s)\,
     n^{1+\ell+(j+\ell)/2+s}
     X^{1+(s+\ell)/2+s}\\
 &\quad\leq C(E,j)C(E,s)\,
               n^{j/2}X^{(5s+4+4\ell)/2}\\
 &\quad\leq C(E,j)C(E,s)\,
               n^{j/2}X^{(9s+8+4\ell)/2}.
 \end{split}
 \tag{PG10}
\]
The last step uses \(X\geq1\).  The weighted sum over pairs
\((L,U)\) equals the left side of (PG2), by orbit--stabilizer
for the finite group \(\operatorname{Aut}(L)\).
Finiteness follows from (MM13), the finitely many possible
discriminants \(b\), and the finite graph count.  Restricting
the lattice collection to \(\mathcal G\) only decreases the
sum, proving the theorem.
\end{proof}

\begin{theorem}[The original normalized primitive-sublattice mean]
\label{thm:normalized-primitive-sublattice-genus-mean}
Write
\[
 \rho=|\operatorname{disc}E|^{1/(2f)}
      =|\operatorname{disc}F|^{1/f},\qquad
 v=\rho/2,\qquad c=1+\log_2(2f).
 \tag{PG11}
\]
Assume \(v\geq1\).  Let \(k\geq3\), \(1\leq s<k\), and let
\(L_0\) be an integral positive definite rank-\(k\) lattice with
\[
 L_0^\#/L_0\simeq\mathcal O_E/T\mathcal O_E
\]
for a rational prime \(T\) splitting completely in \(E\).  Set
\[
 D_0=T^{2f},\qquad
 M_b=(20f\rho^2)^{\,2f(6+c)}.
 \tag{PG12}
\]
For every integer \(n\geq1\),
\[
 \begin{split}
 \sum_{[L]\in\operatorname{gen}(L_0)}
       \frac{|\mathcal U_{k-s,n}(L)|}{|\operatorname{Aut}(L)|}
 \ \leq\ &
 \operatorname{Mass}(L_0)\,
 M_b\rho^f16^{skf}\\
 &{}\cdot(n/D_0)^{k/2}
                   (nD_0)^{5s+5+2c}.
 \end{split}
 \tag{PG13}
\]
Every parameter and exponent on the right is the original one;
in particular \(v\) is the source's inverse-volume parameter.
\end{theorem}

\begin{proof}
\smallskip\noindent\emph{The reference mass in rank one.}
Put \(C_a=C(E,a)\) and \(\ell=\log_2(2f)\).
The lattice \(\mathcal O_E\), with form
\(h(x,y)=\overline x y\), is unimodular.  Thus \(C_1\)
is the mass of its genus.  Its lattices are fractional ideals
\(I\) with \(I\overline I=\mathcal O_E\).

Map their isometry classes to the ideal class group of \(E\).
Fix one ideal \(I_0\) in a nonempty fiber.  A second ideal in
that fiber is \(xI_0\), where \(N_{E/F}x\in\mathcal O_F^\times\).
The assignment
\[
 xI_0\longmapsto
     N_{E/F}x\bmod N_{E/F}(\mathcal O_E^\times)
 \tag{PG14}
\]
is well defined.  If two assignments agree, dividing the
ratio of their generators by an appropriate unit of
\(\mathcal O_E\) gives a norm-one element carrying one ideal
to the other.  Hence the map is injective on isometry classes.
Since the norm subgroup contains
\((\mathcal O_F^\times)^2\), Dirichlet's unit theorem bounds
the fiber by \(2^f\).  Discarding the reciprocal automorphism
weights only increases the sum, and consequently
\[
                             C_1\leq2^fh_E.
 \tag{PG15}
\]

Minkowski's ideal-class bound gives representatives of norm
at most
\[
 (4/\pi)^f\frac{(2f)!}{(2f)^{2f}}
                  \sqrt{|\operatorname{disc}E|}
                       \leq \rho^f.
 \tag{PG16}
\]
For the last inequality, pair the factors \(i\) and
\(2f+1-i\) in \((2f)!\).  This bounds the numerical
prefactor by \((9/(4\pi))^f<1\).  The identity
\(\sqrt{|\operatorname{disc}E|}=\rho^f\) is exact.
The unit theorem and Minkowski bound are used in their
usual number-field forms.\cite[Theorems 4.3 and 5.1]{milne-ant}
The ideal count (MM17) now gives
\[
 h_E\leq(\rho^f)^{1+\ell},\qquad
 C_1\leq2^f\rho^{f(1+\ell)}\leq M_b.
 \tag{PG17}
\]
The last inequality follows directly from \(f\geq1\),
\(\rho\geq2\), and \(M_b=(20f\rho^2)^{2f(7+\ell)}\).
No rank-two class-number bound is needed.

\smallskip\noindent\emph{The comparison across ranks.}
For every integer \(a\geq2\), the Euler factors give
\[
 2^{-f}\leq \zeta_F(a)^{-1}
       \leq L_F(a,\eta^a)\leq\zeta_F(a)\leq2^f.
 \tag{PG18}
\]
Indeed each quadratic-character Euler factor lies between
the corresponding factors of \(\zeta_F(a)^{-1}\) and
\(\zeta_F(a)\); above a rational prime there are at most
\(f\) prime ideals.  Thus \(\zeta_F(a)\leq\zeta(a)^f<2^f\).

Let \(j=k-s\), \(t=\min(j,s)\), and \(u=\max(j,s)\).
Dividing the explicit mass formula (MM1) in ranks \(j,s,k\)
by its rank-one version gives
\[
 \frac{C_jC_s}{C_k C_1}
 =
 \left[
  \rho^{-js-1/2}(2\pi)^{js+1}
  \frac{\prod_{a=1}^{j}(a-1)!\prod_{a=1}^{s}(a-1)!}
       {\prod_{a=1}^{k}(a-1)!}
 \right]^f
 \frac{\prod_{a=2}^{t}L_F(a,\eta^a)}
      {\prod_{a=u+1}^{k}L_F(a,\eta^a)}.
 \tag{PG19}
\]
Empty products equal one.  The factorial ratio is at most
one, since it is
\(\prod_{a=1}^{t}(a-1)!/(u+a-1)!\).
Using (PG18), \(2\pi<8\), \(\rho\geq2\), and \(t\leq s\),
we obtain
\[
 \frac{C_jC_s}{C_k C_1}
       \leq 2^{f(3js+3+2s)}
       \leq \rho^f16^{skf}.
 \tag{PG20}
\]
For the last inequality,
\[
 4sk+1-(3js+3+2s)=s(k+3s-2)-2\geq0
\]
because \(k\geq3\) and \(s\geq1\).
Combining (PG17) and (PG20) proves
\[
                C_jC_s\leq C_kM_b\rho^f16^{skf}.
 \tag{PG21}
\]

\smallskip\noindent\emph{The actual cyclic-discriminant genus.}
Corollary~\ref{cor:cyclic-discriminant-mass} gives
\[
 \operatorname{Mass}(L_0)
   =C_k(1+T+\cdots+T^{k-1})^f
                  \geq C_kD_0^{(k-1)/2}.
 \tag{PG22}
\]
Every lattice in its genus has discriminant \(D_0\).
Apply (PG2) with \(j=k-s\), \(N=D_0\), and then
(PG21)--(PG22).  After extracting
\(\operatorname{Mass}(L_0)M_b\rho^f16^{skf}\), it remains
to bound
\[
 D_0^{-(k-1)/2}n^{(k-s)/2}
                       (nD_0)^{(9s+8+4\ell)/2}
 \leq (n/D_0)^{k/2}(nD_0)^{5s+7+2\ell}.
 \tag{PG23}
\]
The exponent on \(n\) on the right exceeds that on the
left by \(s+3\), and the exponent on \(D_0\) exceeds it
by \((s+5)/2\).  Both bases are at least one.  Since
\(c=1+\ell\), (PG23) is precisely the desired exponent
comparison, and proves (PG13).
\end{proof}

\paragraph{The original downstream requirement.}
The source's primitive-sublattice theorem assumes an
isometry-stable collection of rank-\((j+s)\) integral
lattices of discriminant \(N\) with cyclic dual quotient.
Its right side is
\(\operatorname{Mass}(I_j)\operatorname{Mass}(I_s)
 n^{j/2}(nN)^{(9s+8+4\log_2[E:\mathbf Q])/2}\),
where \(I_j,I_s\) are the specified orthogonal unimodular
lattices.  Their masses equal \(C(E,j),C(E,s)\), respectively.
Thus (PG2) supplies that exact theorem and also covers a
larger class of discriminant modules.  The source's
additional existence hypothesis for orthogonal reference
lattices is not needed in this proof, but its original
reference masses are retained when making the comparison.
\cite{primitive-sublattice-mean-source}

Theorem~\ref{thm:normalized-primitive-sublattice-genus-mean}
also supplies the later normalization by
\(\operatorname{Mass}(L_0)\), including its prescribed field
factor \(M_b\) and all growing-rank exponents.  The fixed-label
pair count still requires the fixed-space isometry estimates
and the original weighted row correspondence.  Neither is
inferred from a count of primitive sublattices alone.
